\chapter{总结与展望}
%=====================================================================
\section{总结}
本文围绕反应堆反电子中微子振荡分析中的若干关键问题展开，目标是在低统计量条件下减少拟合偏差并降低近点实验能量分辨率不足所引入的模型依赖性。研究自“从反应堆出射谱到探测器重构谱”的完整物理链条出发，横跨数值实现、能谱分段（bin/segment）策略与严格统计方法三大层面：一方面通过工程化的性能优化——包括张量化谱生成、稀疏/分块矩阵运算与缓存机制——使得基于海量伪实验的频率学覆盖性检验（如 Feldman–Cousins）在可接受算力下成为可行；另一方面系统评估并量化了重建能谱不同分 bin 策略下的统计偏差与由反应堆模型不确定性引起的系统误差，为联合分析情形下的重建能谱分 bin 选择与反应堆能谱分段提供了可操作的定量依据。

在软件代码框架方面，本文提出并实现了一个基于张量化思想的前向折叠（forward-folded）能谱生成与拟合框架。我们将反应堆中微子能谱、IBD 微分截面、沉积能谱与探测器响应统一表达为矩阵/张量运算链，并利用带状稀疏表示、分块（blocked）构造与结果缓存等技术，显著降低了矩阵乘法与误差函数计算的开销。经基准测试，该框架在常规服务器上相较于未经优化的实现总体加速约一个数量级，同时保持数值精度（相对误差低于 $10^{-6}$）。这一高性能拟合器已作为 JUNO 合作组内部的基础 fitter 被采用，并为后续的大规模 toy-MC、参数扫描与覆盖性研究提供了可行的计算平台。

在分 bin（segment）策略研究方面，本文首先针对 JUNO 重建能谱的分 bin 方案进行了系统评估，量化不同 bin 宽对拟合偏差的影响。研究表明，在低统计量条件下，若分 bin 过细，部分能区内事件数过少，将偏离高斯近似的适用条件，从而在最小化过程中引入可观的参数偏差；随着 bin 宽逐步增大，这种由统计涨落引起的非高斯效应显著减弱。当 bin 宽达到约 100 keV 及以上时，振荡参数的偏差趋于稳定且保持在可接受范围内。据此，本文提出了一套适用于 JUNO 早期运行阶段的重建能谱分 bin 建议。在 JUNO–Daya Bay 联合分析框架下，本文进一步讨论了反应堆能谱的 segment 划分问题。通过大量 Toy-MC 比较固定反应堆模型与包含精细结构扰动的模型在联合拟合中的表现，系统评估了不同段宽下模型依赖引入的附加系统误差。结果表明，当反应堆能谱 segment 宽度取约 300 keV 时，模型相关不确定性对振荡参数的影响最小（小于约 0.3\%），同时避免了因分段过细而超出近点实验分辨能力所产生的非物理结构。该分段策略已被合作组采纳，并体现在 JUNO 首篇物理结果论文的统计方法部分。

在统计方法方面，针对 Wilks 定理在低统计量或非正则情形下的局限性，本文开展了大规模的 Feldman–Cousins（FC）覆盖性检验。结果表明：对太阳区参数 $\Delta m^2_{21}$ 与 $\sin^2\theta_{12}$ 的二维检验，经验 $\Delta\chi^2$ 分布与自由度为 2 的卡方期望一致，可使用 Wilks 近似；但对 $\Delta m^2_{31}$ 的一维检验，经验 1σ 临界值约为 $1.57$，显著高于 Wilks 1$\sigma$对应的临界值1.0，说明在当前低统计量下若直接用 Wilks 阈值会出现覆盖率不足的问题。基于 FC 构造得到的 $\Delta m^2_{31}$ 相对精度约为 $1.4\%$（50 天），并观察到约 8\% 的伪实验会产生两个不相交的置信区间，这进一步强调了在此情形下必须采用基于伪实验的频率学方法以保证置信区间的正确覆盖。

综上所述，本文通过张量化能谱生成与矩阵化重构，构建了一个高性能的反应堆中微子前向折叠拟合框架，使基于海量伪实验的灵敏度评估与严格覆盖性检验在有限计算资源下成为现实；系统提出并验证了合理的重建能谱分 bin 策略与反应堆能谱分段（segment）方案，定量评估了反应堆模型依赖所引入的附加系统误差；同时基于 Feldman–Cousins 方法开展覆盖性研究，为振荡参数置信区间的构造提供了严格的频率学依据，并为 JUNO 合作组早期物理结果的统计处理给出了实质性建议。上述工作不仅提升了振荡参数测量结果的统计可靠性与方法学稳健性，也为未来在更高统计量、更复杂系统不确定性以及联合分析情形下的研究奠定了可推广、可复用的技术与方法基础。


\section{展望}


JUNO 在早期运行中已经展现出令人瞩目的物理潜力。合作组发布的首篇物理结果明确证实了超大体积液体闪烁体探测器结合极高光学性能在中基线探测方案中的可行性。相较于以往全球所有的中微子振荡测量结果，JUNO 成功将测量精度提升了约 1.6 倍，这一里程碑式的进展标志着中微子物理研究正从“发现时代”全面跨越到“高精度测量时代”。未来的工作将在此基础上，利用不断累积的真实物理数据，对探测器的响应函数及能量刻度进行更深层次的校准，以维持并进一步提升这一领先的精度优势。

其次，针对核心物理参数 $\Delta m^2_{31}$ 的精确测量，未来的研究需应对更为严苛的统计学检验。研究表明，即便在累积一年数据量的高统计样本下，$\Delta m^2_{31}$ 的 $\chi^2$ 轮廓仍可能保留显著的非平滑结构，甚至伴随多重局域极小值。这反映出，一年的统计规模仍不足以使实验完全满足 Wilks 定理所需的“大样本渐近极限”。因此依然需要使用Feldman-Cousins方法以保证置信区间的正确覆盖。

与此同时，JUNO 与高分辨率近点探测器 TAO 的联合分析将是实现物理目标的关键突破口。TAO 卓越的能量分辨率为“近点锚定”提供了前所未有的能力，使其能够在极细微的能谱尺度上观测反应堆反中微子的能谱特征。当 TAO 的统计量达到预期水平时，分析精度有望精细到能够有效抑制反应堆谱对小振荡结构的“补偿效应”，从而实现更接近模型独立的（Model-independent）参数测量。

从更宏观的国际物理前沿来看，中微子研究正处于重大发现的前夜。中微子质量顺序（Mass Ordering）的最终判定不仅是 JUNO 的核心科学使命，也是全球中微子实验竞相攻克的制高点。此外，关于中微子是否为马约拉纳费米子的检验——即无中微子双贝塔衰变（$0\nu\beta\beta$）的寻找，以及电弱 CP 破坏相位的精确测定，均是揭示超出标准模型新物理、解释宇宙物质-反物质不对称性的关键钥匙。

综上所述，中微子物理是一个充满生机与活力的研究领域，仍有大量意义重大的未解之谜。随着 JUNO 实验的深入开展及与全球同行的竞争合作，我们有理由相信，在不久的将来，这一领域将产生更多的科学发现，为人类理解微观世界与宇宙演化作出卓越贡献。